\subsection{Welfare comparison}
\label{subsec:rewrite-welfare}

Do the consumption gains from monopoly-financed research compensate for
the restriction on AI use? I evaluate both institutional experiments
against continued competition, using the household's preferences in
equation~\eqref{eq:rewrite-household-problem}. Within each experiment,
I hold the initial state fixed across market structures and research
productivities at each elasticity. Consumption and research after the grant
remain equilibrium choices, not additional initial conditions.

Let $c_t^{\mathrm{mon}}$ and $c_t^{\mathrm{comp}}$ denote consumption per
person under monopoly and continued competition. I report the permanent
proportional change $\lambda$ in competitive consumption that gives the
household the same discounted utility as the monopoly path. With
logarithmic utility and population growth at rate $n$, this
consumption equivalent is
\begin{equation}
 \lambda=\exp\!\left[
 (\rho-n)\int_0^\infty e^{-(\rho-n)t}
 \log\!\left(\frac{c_t^{\mathrm{mon}}}{c_t^{\mathrm{comp}}}\right)\dd t
 \right]-1.
 \label{eq:rewrite-welfare-equivalent}
\end{equation}
Positive values favor monopoly and negative values favor competition.
The comparison includes the transition,
not just the limiting growth rates. Appendix~\ref{app:rewrite-welfare-numerics}
describes the integration and horizon checks.

% Generated by scripts/report_rewrite_welfare.py.
\begin{table}[htbp]
\centering
\caption{Welfare from granting monopoly rights at $\sigma=1.5$}
\label{tab:rewrite-welfare}
\small\setstretch{1.0}
\begin{tabular}{lcrrr}
\toprule
Experiment & $\chi$ & Monopoly growth & Competitive growth & Consumption equivalent\\
\midrule
A & 7.5 & 3.1\% & 1.0\% & 130.3\%\\
 & 1.5 & 3.1\% & 1.0\% & 11.1\%\\
\midrule
B & 7.5 & 1.0\% & 5.7\% & -60.7\%\\
 & 1.5 & 1.0\% & 5.7\% & -65.0\%\\
\bottomrule
\end{tabular}
\par\vspace{4pt}
\parbox{0.95\textwidth}{\footnotesize\emph{Notes:} Growth rates refer to long-run consumption per person. Within each experiment, all paths inherit the same competitive initial state. In A, monopoly enables AI-dominated growth, as in Section~\ref{subsec:rewrite-competitive-transition}. In B, monopoly eliminates it, as in the supplementary growth-reversal experiment, with different initial stocks and a lower efficiency bound. The consumption equivalent is the permanent proportional change in competitive consumption that gives the representative household the same discounted utility as the monopoly path. Positive values favor monopoly and negative values favor competition. Preferences are unchanged across both experiments.}
\end{table}


Experiment A in Table~\ref{tab:rewrite-welfare} extends the transition in
Section~\ref{subsec:rewrite-competitive-transition} to all four elasticities
in the supplementary simulations. At $\sigma=0.9$, monopoly reduces
welfare under both research productivities, with consumption equivalents
of $-29.8\%$ and $-31.9\%$. At $\sigma=1$ and $1.1$, monopoly raises
welfare with $\chi=7.5$ but lowers it with $\chi=1.5$. In all six cases,
both market structures approach consumption growth of $1.0\%$.
Research can therefore compensate for monopoly pricing without raising
long-run growth, but whether it does so depends on the consumption levels
and transition generated by research productivity.

At $\sigma=1.5$, the same experiment shows consumption-equivalent gains of
$130.3\%$ with $\chi=7.5$ and $11.1\%$ with $\chi=1.5$.
Both monopoly economies approach consumption-per-person growth of
$3.1\%$, compared with $1.0\%$ under continued competition. Their identical
limiting growth rates therefore conceal substantial differences in
welfare. Earlier efficiency gains under higher research productivity
raise consumption sooner and leave it on a higher subsequent path.

Experiment B evaluates the opposite outcome in the
\href{https://oliverpardo1979.github.io/ai-growth-and-labor/\#subsec:rewrite-monopoly-growth-reversal}
{supplementary growth-reversal experiment}. Monopoly reduces long-run
consumption growth from $5.7\%$ under competition to $1.0\%$.
The consumption equivalents are $-60.7\%$ with $\chi=7.5$ and
$-65.0\%$ with $\chi=1.5$. The household would be indifferent between
each monopoly path and a permanent reduction of that magnitude in
competitive consumption. Lower compute use and investment initially
release resources for consumption, but these early gains do not
compensate for the subsequent consumption losses.

These comparisons concern the representative household, which owns capital and
the developer as well as supplying labor. The initial wage losses do not
by themselves determine its welfare. These rankings apply to
the specified illustrative experiments,
not to all parameter values or initial states, and do not establish
that granting monopoly rights is an optimal policy.
